%==============================================================================
% MICRO submission draft.
%
% FORMAT: MICRO has used the IEEE conference format for many years, which is
% what this file targets. Some recent architecture conferences (ISCA, ASPLOS)
% have moved to ACM acmart/sigconf, and MICRO's template has varied by year --
% CHECK THE CURRENT CFP before submitting. To switch, replace the two lines
% below with:
%     \documentclass[sigconf,nonacm]{acmart}
% and delete the \IEEEoverridecommandlockouts / \thanks block.
%
% Neither IEEEtran.cls nor acmart.cls was installed on the machine this was
% written on, so THIS FILE HAS NOT BEEN COMPILED. Install with e.g.
%     sudo apt install texlive-publishers texlive-latex-extra  # IEEEtran
%     tlmgr install acmart                                     # acmart
%==============================================================================
\documentclass[conference]{IEEEtran}
\IEEEoverridecommandlockouts

\usepackage{booktabs}
\usepackage{amsmath}
\usepackage{graphicx}
\usepackage{xcolor}
\usepackage{siunitx}
\usepackage[hidelinks]{hyperref}

\sisetup{group-separator={,}, group-minimum-digits=4}

\newcommand{\cfgb}{\textsf{(b)}}
\newcommand{\cfgc}{\textsf{(c)}}
\newcommand{\cfgd}{\textsf{(d)}}

\begin{document}

\title{Absorbing the Transient KV Write Stream:\\
A Phase-Change Memory Tier Between HBM and High-Bandwidth Flash}

\author{\IEEEauthorblockN{Rishabh Chamoli}
\IEEEauthorblockA{Indian Institute of Technology Guwahati\\
rishabh.chamoli@iitg.ac.in}}

\maketitle

%==============================================================================
\begin{abstract}
%==============================================================================
High-Bandwidth Flash (HBF) promises roughly $16\times$ the capacity of HBM at
comparable bandwidth, making it an attractive destination for the KV cache of
large-language-model inference. Proposals such as H3 place HBF behind HBM in a
daisy-chained hierarchy, but confine it to \emph{read-only} data---model weights
and shared, pre-computed KV---precisely because NAND write endurance is the
weak point of the technology. We ask what happens to that hierarchy under
general write-heavy serving, where the flash tier does receive a transient write
stream, and whether a small bit-alterable NVM tier can absorb it.

We add a \SI{43.4}{\gibi\byte}-per-GPU phase-change memory (PCM) tier,
physically stackable on top of each HBF stack within its existing footprint, and
evaluate it in a discrete-event LLM serving simulator against the Mooncake
FAST'25 traces on an 8-GPU LLaMA2-70B-GQA cluster. At the buildable capacity, a
write-back admission policy cuts flash write volume by \SI{51.2}{\percent}
(conversational) and \SI{52.6}{\percent} (tool/agent); adding a read-gated
demotion policy reaches \SI{98.8}{\percent} and \SI{98.3}{\percent}, a
$90.3\times$ and $86.4\times$ work-normalized lifetime improvement.

Three results temper this. First, flash wear is governed by tier capacity
against workload footprint, \emph{not} by batch size: across a $12\times$
concurrency range driving peak HBM occupancy from \SI{25}{\percent} to
\SI{99.2}{\percent}, flash bytes per request move by under
\SI{1.1}{\percent}. Second, \SI{43.4}{\gibi\byte} per GPU sits at
\SIrange{5.3}{10.2}{\percent} of the workload footprint, on the flat foot of the
capacity curve; the knee at which flash goes fully read-only is $19\times$
further out. Third, the read-gated policy buys its write reduction at a
\SIrange{9.6}{40.7}{\percent} cost in prefix-cache hit rate. We additionally
report a pre-registered falsification test of the central mechanism, which the
architecture passes.
\end{abstract}

\begin{IEEEkeywords}
LLM inference, KV cache, high-bandwidth flash, phase-change memory, write
endurance, memory hierarchy.
\end{IEEEkeywords}

%==============================================================================
\section{Introduction}
%==============================================================================
The key--value (KV) cache of a transformer dominates the memory budget of
large-language-model (LLM) serving. At LLaMA2-70B with grouped-query attention,
tensor-parallel degree two, and FP16 precision, a single token costs
\SI{163840}{\byte} of KV state; a \num{13000}-token prompt therefore occupies
roughly \SI{2}{\gibi\byte} before a single token is generated. HBM capacity, not
compute, is what bounds concurrency.

High-Bandwidth Flash (HBF) is a recent response: a NAND stack built to HBM-like
bandwidth and placed in the same package, offering on the order of $16\times$
the capacity. The H3 proposal~\cite{h3} daisy-chains HBF behind HBM, routing
through an address decoder in the HBM base die and hiding flash's microsecond
read latency behind a \SI{40}{\mega\byte} SRAM prefetch buffer. Crucially, H3
confines HBF to \emph{read-only} data---weights and shared, pre-computed KV---and
argues on that basis that NAND's limited write endurance is not a disadvantage.

That argument does not survive contact with general serving. When the flash tier
backs a KV store that absorbs spill from HBM, it receives a sustained write
stream, and TLC NAND rated at \num{3000} program/erase cycles becomes the
limiting component of the package. This paper asks whether inserting a small
amount of \emph{bit-alterable} non-volatile memory---phase-change memory---between
HBM and HBF absorbs enough of that stream to make the hierarchy viable, and at
what cost.

\paragraph{Contributions}
\begin{itemize}
  \item A concrete PCM tier derived from published die parameters
        (\SI{43.4}{\gibi\byte} per GPU) that stacks vertically on the existing
        HBF stacks and consumes no HBF capacity (\S\ref{sec:arch}).
  \item A measurement of what write-absorption policy buys at that capacity:
        \SIrange{51}{53}{\percent} write reduction from write-back admission
        alone, and \SIrange{98}{99}{\percent} with read-gated demotion
        (\S\ref{sec:headline}).
  \item The negative result that flash wear is invariant to batch size and HBM
        occupancy, which removes an entire class of proposed deployment rules
        (\S\ref{sec:batch}).
  \item A capacity sweep locating the knee at $19\times$ the buildable budget,
        establishing that the design point is defensible as \emph{buildable} but
        is not a capacity optimum (\S\ref{sec:capacity}).
  \item A pre-registered falsification test of the mechanism by which HBM
        pressure is assumed to reach the flash tier (\S\ref{sec:falsifier}).
  \item An explicit statement of what the measurements do \emph{not} cover
        (\S\ref{sec:limits}), including the fact that generated KV is bounded at
        \SIrange{1.9}{2.8}{\percent} of tokens on these traces.
\end{itemize}

%==============================================================================
\section{Background and Motivation}
%==============================================================================

\subsection{KV cache offload}
Production serving systems already treat the KV cache as a storage hierarchy.
vLLM's paged attention~\cite{vllm} manages HBM in fixed-size blocks and reuses
shared prefixes across requests; Mooncake~\cite{mooncake} extends this to a
disaggregated store with a DRAM memory tier and an SSD offload tier, admitting
blocks on prefill and evicting by LRU. The economics are attractive because
prompts are heavily shared: on the traces used here, \num{4080885} prompt-block
instances deduplicate to \num{2695220} distinct blocks, a $1.51\times$ factor.

\subsection{H3: HBM plus high-bandwidth flash}
H3~\cite{h3} assumes a B200-class GPU with \SI{192}{\giga\byte} of HBM3e across
eight cubes at \SI{8}{\tera\byte\per\second}, and adds eight HBF stacks of
\SI{384}{\giga\byte} each---\SI{3}{\tera\byte} at comparable aggregate
bandwidth. Two structural points are often modelled incorrectly and matter here:

\begin{enumerate}
  \item \textbf{No HBM is surrendered.} The HBM cubes remain on the GPU
        shoreline; HBF is \emph{daisy-chained} behind them, with an address
        decoder and router in the HBM base die splitting traffic across one
        unified address space. There is no per-stack bandwidth derate. A
        side-by-side layout that trades HBM capacity for flash models a
        different machine.
  \item \textbf{Latency is hidden, not low.} A \SI{40}{\mega\byte} SRAM latency
        hiding buffer occupying \SI{8.06}{\milli\meter\squared} of a
        \SI{121}{\milli\meter\squared} base die (\SI{6.7}{\percent}) prefetches
        from HBF. The mechanism is \emph{predictability}: layer $N{+}1$'s
        weights are known while layer $N$ computes. It is a read-side buffer; it
        hides read latency, not write latency, and it creates no bandwidth.
\end{enumerate}

\subsection{Why endurance becomes the problem}
H3 targets read-only data and argues that low write endurance is therefore not a
disadvantage. We target \emph{generated and transient} KV on the same hardware.
Under that workload HBF receives the spill stream from a KV store, and TLC NAND
at \num{3000} P/E cycles must absorb it. Taking the cluster flash budget of
\SI{24}{\tebi\byte} at \num{3000} P/E and assuming a write-amplification factor
of one with ideal wear levelling gives
$24 \times 2^{40} \times 3000 = \SI{7.92e16}{\byte}$ of lifetime writes---a
figure we use only as a \emph{ratio} denominator, never to quote a lifetime in
days, because the far tier has no shared-bandwidth model in this work.

PCM is the natural absorber: it is bit-alterable, so it does not pay NAND's
erase-block penalty, and its write endurance is orders of magnitude higher.
Its cost per bit and its exact endurance figure are the open variables.

%==============================================================================
\section{Architecture}\label{sec:arch}
%==============================================================================

\begin{table}[t]
\centering
\caption{Per-GPU memory hierarchy. HBM and HBF figures follow
H3~\cite{h3} for a B200-class part; the PCM tier is this work's addition.}
\label{tab:hierarchy}
\small
\begin{tabular}{llll}
\toprule
\textbf{Tier} & \textbf{Capacity} & \textbf{Bandwidth} & \textbf{Role} \\
\midrule
HBM3e & \SI{192}{\giga\byte} & \SI{8}{\tera\byte\per\second} & memory tier \\
      & 8 cubes $\times$ \SI{24}{\giga\byte} & \SI{1}{\tera\byte\per\second}/cube & \emph{all} KV \\[2pt]
PCM   & \textbf{\SI{43.4}{\gibi\byte}} & --- & write absorber \\
      & 8 stacks $\times$ \SI{5.43}{\gibi\byte} & & \\[2pt]
HBF   & \SI{3}{\tera\byte} & \SI{8}{\tera\byte\per\second} & offload tier \\
      & 8 stacks $\times$ \SI{384}{\giga\byte} & & \\
\bottomrule
\end{tabular}
\end{table}

\subsection{Placement and capacity of the PCM tier}
We place two PCM dies on top of each HBF stack, \emph{within the existing stack
footprint}. PCM does not displace NAND dies laterally and therefore does not
reduce HBF capacity. At \SI{35}{\milli\meter\squared} per die and a density of
\SI{0.62}{\giga\bit\per\milli\meter\squared}:
\begin{align}
2 \times 35 \times 0.62 &= \SI{43.4}{\giga\bit} \text{ per stack} \nonumber\\
                        &= \SI{5.43}{\gibi\byte} \text{ per stack} \nonumber\\
8 \times \SI{5.43}{\gibi\byte} &= \mathbf{\SI{43.4}{\gibi\byte}} \text{ per GPU.}
\end{align}
For comparison, the NAND side is TLC at
\SI{29}{\giga\bit\per\milli\meter\squared}, so a \SI{1}{\tera\bit} die occupies
roughly \SI{35.3}{\milli\meter\squared} and holds \SI{128}{\gibi\byte}. The PCM
tier is thus about \SI{1.4}{\percent} of the flash capacity it protects.

\subsection{Write-absorption policies}
The store exposes two orthogonal policy knobs, giving the three configurations
evaluated throughout (Table~\ref{tab:configs}):

\begin{itemize}
  \item \textbf{Admission write policy.} Under \emph{write-through}, every
        admission to the memory tier immediately writes the block through to
        flash; a later demotion is then a byte-identical duplicate, and no
        eviction-time policy can remove more than that duplicate half. Under
        \emph{write-back}, admission spends no flash write at all, and a block
        reaches flash only if demotion later decides to send it.
  \item \textbf{Demotion policy.} Under \emph{always}, every memory-tier
        eviction demotes its victim. Under \emph{if\_read}, only victims that
        were read back at least once while resident are demoted; never-read
        victims are dropped, on the grounds that the bytes bought no cache hit.
\end{itemize}

\begin{table}[t]
\centering
\caption{Evaluated configurations.}
\label{tab:configs}
\small
\begin{tabular}{llll}
\toprule
 & \textbf{Memory tier} & \textbf{Admission} & \textbf{Demotion} \\
\midrule
\cfgb\ HBF only      & DRAM & write-through & always \\
\cfgc\ +PCM          & PCM  & write-back    & always \\
\cfgd\ +PCM, gated   & PCM  & write-back    & if\_read \\
\bottomrule
\end{tabular}
\end{table}

Configuration \cfgb\ is the HBF-only baseline: it is what the H3 hierarchy does
if a KV store is simply pointed at it. Writing it this way makes the comparison
a policy comparison at fixed hardware, which is the honest form of the question.

%==============================================================================
\section{Methodology}\label{sec:method}
%==============================================================================

\subsection{Simulator}
We extend TokenSim~\cite{tokensim}, a discrete-event simulator for LLM inference
systems that models request arrival, prefill/decode scheduling, paged KV
management, distributed execution and roofline-derived latency. Its Mooncake
store is modelled as a memory tier plus an offload tier with LRU eviction and a
per-operation latency function
$t = t_\text{fixed} + \text{bytes} / \text{BW}$.

The latency hiding buffer is not simulated as a buffer; its effect is folded
into the flash read parameters as \SI{0.1}{\micro\second} fixed latency at
\SI{8}{\tera\byte\per\second}. This is a \emph{perfect-prefetch upper bound},
not a measurement, and it is deliberately generous to the HBF-only baseline---it
makes flash reads cheap, which is conservative for the argument we are making.
Flash writes are NAND-realistic at \SI{200}{\micro\second} and
\SI{3}{\giga\byte\per\second} and are \emph{not} hidden, since the buffer is a
read-side prefetcher.

\subsection{Platform and capacity bookkeeping}
The cluster is eight B200-class workers running LLaMA2-70B-GQA
($N_\text{head}=64$, 8 KV heads, $N_\text{layer}=80$, $d_\text{model}=8192$) at
tensor-parallel degree 2 and data-parallel degree 4. KV precision is FP16, which
is hardcoded in the harness; FP8 would halve KV bytes per token and improve every
write-volume result, so FP16 \emph{understates} the benefit.

This gives \SI{163840}{\byte} per token and a \num{16}-token block of
\SI{2621440}{\byte}. The store is a single cluster-wide pool, so capacities are
cluster totals:
\begin{align}
\text{PCM pool} &= \SI{43.4}{\gibi\byte} \times 8 = \SI{347.2}{\gibi\byte} = \num{142213}~\text{blk} \nonumber\\
\text{HBF pool} &= \SI{3}{\tebi\byte} \times 8 = \SI{24}{\tebi\byte} = \num{10066329}~\text{blk} \nonumber\\
\text{HBM/rank} &= \SI{192}{\gibi\byte} = \num{78643}~\text{blk.} \nonumber
\end{align}

\subsection{Workload}
We use the Mooncake FAST'25 conversational and tool/agent
traces~\cite{mooncake}, first \num{5000} requests each, replayed at an offered
load of \SI{30}{\per\second}. Table~\ref{tab:workload} gives the footprints.
These traces are a \textbf{stand-in}: the target workload is cache-augmented
generation, whose parameters are not yet fixed, and every result should be read
with that substitution in mind.

\begin{table}[t]
\centering
\caption{Workload footprint, $W=\num{5000}$ requests. Distinct blocks are
counted with the store's own truncation.}
\label{tab:workload}
\small
\begin{tabular}{lrr}
\toprule
 & \textbf{Conversational} & \textbf{Tool/Agent} \\
\midrule
prompt-block instances & \num{4080885} & \num{2908548} \\
\textbf{distinct blocks} & \textbf{\num{2695220}} & \textbf{\num{1397479}} \\
as bytes & \SI{6580}{\gibi\byte} & \SI{3412}{\gibi\byte} \\
deduplication & $1.51\times$ & $2.08\times$ \\
input tokens, mean & \num{13066} & \num{9315} \\
output tokens, mean & \num{346} & \num{185} \\
\bottomrule
\end{tabular}
\end{table}

\subsection{Metrics and reproducibility}
All results are \textbf{work-normalized}: bytes per request, bytes per token, or
ratios. We quote no absolute lifetimes in days, because the far tier has no
shared-bandwidth model. Every run reports the occupancy and capacity
\emph{achieved}, not the knob value set, and the achieved batch per scheduler is
reported alongside, so that a configuration which changes throughput cannot
silently change the operating point.

Each experimental campaign is gated on reproducing a previously recorded
configuration to every digit before any new number is trusted. The workload
rebuild used here reproduces twelve independently recorded values exactly---block
instance counts, distinct-block counts, input/output length distributions, and
five simulator outputs per trace---on the first attempt.

%==============================================================================
\section{Evaluation}
%==============================================================================

\subsection{Write absorption at the design point}\label{sec:headline}

\begin{table}[t]
\centering
\caption{HBM-saturated operation, $W=\num{5000}$, \SI{30}{\per\second} offered,
peak HBM occupancy \SIrange{99.1}{99.2}{\percent} in all six runs.
\emph{B/req} is HBF write bytes per request.}
\label{tab:headline}
\small
\begin{tabular}{lrrr}
\toprule
\textbf{Conversational} & \cfgb & \cfgc & \cfgd \\
\midrule
HBF write blocks      & \num{6461439} & \num{3150432} & \textbf{\num{80968}} \\
HBF B/req ($\times 10^9$) & 3.388 & 1.652 & \textbf{0.042} \\
blocks demoted / dropped & \num{3159613} / 0 & \num{3150432} / 0 & \num{80968} / \num{3563313} \\
prefix hit rate       & 0.1631 & 0.1641 & 0.0942 \\
achieved batch/sched  & 43.61 & 42.90 & 41.55 \\
\textbf{write reduction vs \cfgb} & --- & \textbf{\SI{51.24}{\percent}} & \textbf{\SI{98.75}{\percent}} \\
\midrule
\textbf{Tool/Agent} & \cfgb & \cfgc & \cfgd \\
\midrule
HBF write blocks      & \num{3234691} & \num{1534945} & \textbf{\num{55703}} \\
HBF B/req ($\times 10^9$) & 1.696 & 0.805 & \textbf{0.029} \\
blocks demoted / dropped & \num{1546239} / 0 & \num{1534945} / 0 & \num{55703} / \num{1694505} \\
prefix hit rate       & 0.4014 & 0.4033 & 0.3733 \\
achieved batch/sched  & 48.44 & 47.32 & 44.72 \\
\textbf{write reduction vs \cfgb} & --- & \textbf{\SI{52.55}{\percent}} & \textbf{\SI{98.28}{\percent}} \\
\bottomrule
\end{tabular}
\end{table}

Table~\ref{tab:headline} gives the headline. Moving admission from write-through
to write-back---configuration \cfgb\ to \cfgc---halves flash write volume, and it
does so for a structural reason: under write-through every admission already
spent a flash write, so the demotion that follows is a byte-identical duplicate.
Write-back removes that duplicate half exactly. The residual is the genuine
spill.

Adding read-gated demotion (\cfgd) removes almost all of what remains, reaching
\SI{98.75}{\percent} and \SI{98.28}{\percent} reduction. Expressed as
work-normalized lifetime against the flash endurance budget, that is
$\mathbf{90.3\times}$ and $\mathbf{86.4\times}$ more requests served before
wear-out. Generated token counts are identical across configurations within each
trace, so the per-request and per-token multipliers coincide exactly.

Achieved batch is matched to within \SI{4.7}{\percent} and \SI{7.7}{\percent}
across the three configurations, preemption and recomputation counts are zero
throughout, and the conservation identity
$\text{demoted} + \text{dropped} = \text{memory evictions}$ holds in every run.

\subsection{Wear is invariant to concurrency}\label{sec:batch}

A natural deployment rule would be ``below concurrency $N$, the flash tier stays
effectively read-only.'' It does not exist. Table~\ref{tab:batch} sweeps the
per-scheduler concurrency cap over a $12\times$ range.

\begin{table}[t]
\centering
\caption{Concurrency sweep, configuration \cfgb, conversational trace.
HBF bytes per request are flat.}
\label{tab:batch}
\small
\begin{tabular}{rrrr}
\toprule
\textbf{cap} & \textbf{achieved batch} & \textbf{peak HBM} & \textbf{HBF B/req} \\
\midrule
8   & 3.94  & \SI{25.0}{\percent} & \num{3377115693} \\
16  & 7.83  & \SI{39.8}{\percent} & \num{3370987815} \\
32  & 15.46 & \SI{62.1}{\percent} & \num{3383093625} \\
48  & 23.17 & \SI{85.7}{\percent} & \num{3370675339} \\
64  & 30.81 & \SI{99.0}{\percent} & \num{3376704651} \\
96  & 42.16 & \SI{99.1}{\percent} & \num{3385003082} \\
192 & 43.61 & \SI{99.2}{\percent} & \num{3387654930} \\
256 & 43.61 & \SI{99.2}{\percent} & \num{3387654930} \\
\midrule
\multicolumn{3}{l}{\textbf{spread}} & \textbf{\SI{0.50}{\percent}} \\
\bottomrule
\end{tabular}
\end{table}

Over a $12\times$ concurrency range that drives peak HBM occupancy from
\SI{25}{\percent} to \SI{99.2}{\percent}, flash bytes per request move by
\SI{0.50}{\percent} (conversational) and \SI{1.10}{\percent} (tool/agent)---noise,
and not even monotone. The PCM multiplier is flat at $2.05\times$ and
$2.09\times$ at every point.

The reason is structural. Blocks enter the store from the prefill admission
path, once per request, and the write stream is therefore set by the workload's
block footprint against the pool capacity, not by how many requests are resident
simultaneously. HBM occupancy changes how \emph{long} a block sits before
eviction, not how \emph{many} blocks are evicted.

The practical consequence is that every endurance figure in this paper is
\emph{concurrency-invariant}, which makes it a property of the policy and the
pool size rather than of one operating point---considerably more robust than a
figure measured at a single batch.

\subsection{Capacity: the design point is buildable, not optimal}\label{sec:capacity}

\begin{table}[t]
\centering
\caption{PCM capacity sweep, conversational trace (footprint
\num{2695220} blocks). ``Requests'' is work-normalized requests to wear-out
at WAF${=}1$.}
\label{tab:capacity}
\small
\begin{tabular}{rrrrr}
\toprule
\textbf{GiB/GPU} & \textbf{\% footprint} & \textbf{\cfgb\ B/req} & \textbf{\cfgc\ B/req} & \textbf{\cfgb$\to$\cfgc} \\
\midrule
8.2   & \SI{1}{\percent}   & \num{3443815088} & \num{1709646021} & $2.01\times$ \\
32.9  & \SI{4}{\percent}   & \num{3407126987} & \num{1670769541} & $2.04\times$ \\
\textbf{43.4} & \textbf{\SI{5.3}{\percent}} & \textbf{\num{3387654930}} & \textbf{\num{1651733692}} & $\mathbf{2.05\times}$ \\
82.3  & \SI{10}{\percent}  & \num{3312823828} & \num{1582069973} & $2.09\times$ \\
205.6 & \SI{25}{\percent}  & \num{3097571623} & \num{1369868075} & $2.26\times$ \\
411.3 & \SI{50}{\percent}  & \num{2744303747} & \num{1016600199} & $2.70\times$ \\
822.5 & \SI{100}{\percent} & \num{2037767995} & \num{310064447}  & $6.57\times$ \\
1233.8 & \SI{150}{\percent} & \num{1725419749} & \textbf{0} & $\infty$ \\
\bottomrule
\end{tabular}
\end{table}

Table~\ref{tab:capacity} sweeps PCM capacity across \num{150} fold. Three things
follow.

\textbf{The buildable point is on the flat foot of the curve.}
\SI{43.4}{\gibi\byte} per GPU is \SI{5.3}{\percent} of the conversational
footprint and \SI{10.2}{\percent} of the tool/agent one. Below roughly half the
footprint, PCM capacity is nearly irrelevant to flash lifetime: the $2.05\times$
benefit at the design point is bought by the \emph{write-back policy}, not by the
capacity. This matters for how the tier should be justified---it is defensible as
the buildable point, and the benefit is real, but it must not be presented as a
capacity optimum.

\textbf{The knee exists and is unbuildable.} Flash write volume collapses to zero
only when the pool can hold the entire workload footprint, at roughly
\SI{822}{\gibi\byte} per GPU---$19\times$ the budget. Only write-back can reach
zero at all; write-through has a hard floor of one flash write per distinct
block at any capacity.

\textbf{Residency and write volume answer different questions.} Mean PCM
residency crosses the median time-to-first-read at $1.8\times$ to
$2.8\times$ the design point, i.e.\ \SIrange{77}{120}{\gibi\byte} per GPU. The tier is within a factor of two to three of being large
enough to \emph{hold} a block until its reuse arrives, while being $19\times$
short of keeping it off flash entirely.

\textbf{Capacity buys writes and latency, never hit rate.} Prefix-cache hit rate
is flat at \numrange{0.1624}{0.1646} and \numrange{0.4007}{0.4036} across the
whole \num{150}-fold sweep. A block demoted to flash remains in the store index
and is still found by lookup---it is merely served from the slower tier. PCM is a
write absorber and, on this evidence, \emph{only} a write absorber.

\subsection{The cost of read-gated demotion}

Configuration \cfgd's \SI{98}{\percent} write reduction is not free. Dropping
never-read victims removes blocks that a later request would have hit, so the
prefix-cache hit rate falls by \SI{40.7}{\percent} (conversational,
$0.1643 \to 0.0974$) and \SI{9.6}{\percent} (tool/agent,
$0.4037 \to 0.3650$) at matched batch, with a corresponding rise in effective
prefill tokens. The two traces differ sharply because tool/agent has far more
in-window reuse, so fewer of its victims are never-read.

This is the central engineering trade of the design, and it is workload
dependent in a way that deserves emphasis: under a trace variant that perturbs
only \SIrange{2}{5}{\percent} of the footprint, the same cost moves to
\SI{27.7}{\percent} and \SI{3.3}{\percent}. It should be quoted with that
sensitivity attached, never as a constant.

\subsection{A falsification test}\label{sec:falsifier}

The flat concurrency curve of \S\ref{sec:batch} invites an objection: perhaps it
is an artefact of the simulator, in which HBM eviction does not feed the store.
If an HBM eviction \emph{did} write its victim to the store, flash bytes per
request should depend on HBM occupancy.

We implemented that path behind a flag and fixed the decision rule in advance:
with the hook enabled, if flash bytes per request at the high-batch point differ
from the low-batch point by less than \SI{5}{\percent} on both traces, the
hypothesis is dead. The result is \SI{+0.00}{\percent} (conversational) and
\SI{-0.07}{\percent} (tool/agent)---the hypothesis is dead.

The mechanism is now measurable. Eviction \emph{volume} itself barely moves with
batch (\SI{+1.66}{\percent} and \SI{+1.42}{\percent} over the same $11\times$
range), because the block allocator hands out never-touched blocks before
recycling cached ones, so the eviction rate tracks total allocations rather than
instantaneous occupancy. Moreover \SIrange{97.4}{98.0}{\percent} of evicted
blocks carry a key the store already holds, making their write a no-op refresh.
Enabling the hook raises flash writes by only
\SIrange{1.85}{2.85}{\percent}.

We report this because it is the kind of result that is easy not to look for:
the flat curve is the paper's most load-bearing negative claim, and it survives
a direct attempt to break it.

%==============================================================================
\section{Limitations}\label{sec:limits}
%==============================================================================
We state these in full because they bound every number above.

\textbf{The store holds prompt blocks only.} Generated KV has no key space in
the simulator and therefore never enters the store. The defensible claim is
narrower than the motivation: \emph{PCM as a write-back buffer in front of HBF
reduces prefix-cache offload writes by \SI{51}{\percent} (always) or
\SI{99}{\percent} (if\_read)}. Closing this gap requires a key space for decode
blocks. We bound what it could add: generated tokens are
\SI{2.58}{\percent} and \SI{1.94}{\percent} of all tokens on these traces
(\SIrange{4.0}{4.1}{\percent} measured against the store's deduplicated write
population), so on \emph{this} workload the ceiling is a few percent. These are
prompt-heavy traces; a workload with long generations would move that bound.

\textbf{The workload is a substitute.} The target is cache-augmented generation;
its corpus size, sharing factor and per-request transient volume are not yet
fixed, and the Mooncake traces stand in for it.

\textbf{FP16 is a harness limitation.} KV precision is hardcoded. FP8 would halve
KV bytes per token, doubling effective tier capacity in blocks and improving
every write-volume and endurance figure. Reporting at FP16 is the conservative
choice.

\textbf{Flash reads are modelled generously.} The perfect-prefetch read
parameters flatter the HBF-only baseline, and the prefetch argument that
justifies them applies to H3's predictable, read-only accesses---not to the
demand-driven hits on transient KV that this workload generates.

\textbf{Lifetime \emph{ratios} are fragile; bytes per request are not.} Because
\cfgd's flash writes are near zero, small absolute changes move the ratio by
large factors: under the trace variant mentioned above, $80.4\times$ becomes
$230.3\times$. Work-normalized bytes per request is the robust figure.

\textbf{PCM endurance is unpriced.} We have no sourced endurance figure for the
PCM tier, so the write traffic it absorbs is reported but its own wear is not
evaluated. This is the single largest open parameter in the design.

\textbf{No shared-bandwidth model for the far tier.} Hence no lifetime in days
anywhere in this paper.

Separately, the trace's \num{512}-token hash granularity against the simulator's
\num{16}-token blocks makes a continuation unable to reuse its predecessor's
prompt tail in \SIrange{97.9}{98.7}{\percent} of cases, affecting
\SIrange{1.35}{1.95}{\percent} of distinct blocks. Reported hit rates are
therefore floors. Comparisons between configurations at a fixed workload are
unaffected, since both arms see the same trace.

%==============================================================================
\section{Related Work}
%==============================================================================
KV cache management in HBM is well studied: paged attention~\cite{vllm}
eliminated fragmentation and enabled prefix sharing, and
Mooncake~\cite{mooncake} extended it into a disaggregated store with DRAM and
SSD tiers. Our work treats that store's offload tier as the object of study
rather than its mechanism.

On the hardware side, H3~\cite{h3} is the direct antecedent and supplies our
topology, capacities and latency-hiding argument. We depart from it in workload
rather than structure: H3 confines flash to read-only data and argues endurance
is therefore a non-issue; we accept the structure and test the write-heavy case
it excludes. Flash-level write amplification and garbage collection behaviour,
which would refine the WAF${=}1$ assumption used here, are addressable with an
SSD-level simulator such as MQSim~\cite{mqsim}; that analysis is future work.

Phase-change memory as a DRAM-adjacent tier has a long literature; what is new
here is the specific role---a thin write absorber protecting a NAND tier inside a
GPU package---and the finding that its benefit at buildable capacity comes from
policy rather than from capacity.

%==============================================================================
\section{Conclusion}
%==============================================================================
A \SI{43.4}{\gibi\byte}-per-GPU PCM tier, stackable within the existing HBF
footprint, absorbs \SIrange{51}{53}{\percent} of the flash write stream through
write-back admission alone, and \SIrange{98}{99}{\percent} with read-gated
demotion, for a $90\times$ work-normalized lifetime improvement. The benefit at
that capacity is bought by policy, not by capacity: the design point sits on the
flat foot of the curve, $19\times$ short of the knee.

Two negative results are, we think, as useful as the positive one. Flash wear is
invariant to batch size and HBM occupancy, which removes a class of deployment
rules that would otherwise look plausible; and that invariance survives a direct
falsification attempt in which HBM evictions are wired into the store. The
remaining gap is the one we name rather than paper over: generated KV does not
yet enter the store, and on these prompt-heavy traces it is bounded at a few
percent of all KV. Extending the key space to decode blocks, and pricing PCM's
own endurance, are the next steps.

%==============================================================================
% NOTE: bibliographic details below are INCOMPLETE where marked. Verify and
% complete them before submission -- do not submit the placeholders.
%==============================================================================
\begin{thebibliography}{9}

\bibitem{h3}
SK hynix,
``H3: a hybrid HBM and high-bandwidth flash memory hierarchy,''
\emph{IEEE Computer Architecture Letters}, 2026.
% TODO: complete authors, volume, issue, page range, DOI.

\bibitem{tokensim}
``TokenSim: enabling hardware and software exploration for large language model
inference systems,'' arXiv:2503.08415, 2025.
% TODO: complete author list.

\bibitem{mooncake}
``Mooncake: a KVCache-centric disaggregated architecture for LLM serving,''
in \emph{Proc. USENIX FAST}, 2025.
% TODO: complete author list and page range.

\bibitem{vllm}
W. Kwon \emph{et al.},
``Efficient memory management for large language model serving with
PagedAttention,''
in \emph{Proc. ACM SOSP}, 2023.

\bibitem{mqsim}
A. Tavakkol, J. G\'{o}mez-Luna, M. Sadrosadati, S. Ghose, and O. Mutlu,
``MQSim: a framework for enabling realistic studies of modern multi-queue SSD
devices,''
in \emph{Proc. USENIX FAST}, 2018.

\end{thebibliography}

\end{document}
